\documentclass[../../main.tex]{subfiles}
\begin{document}
% SOURCE p=230; printed=215; object=ch07-positive-exercises-title; region=middle-lower
\Needspace{12\baselineskip}
\subsection*{习题}
% SOURCE p=230; printed=215; object=ch07-ex01; number=1; region=middle-lower
\begin{exercise}\label{exercise:ch07-closed-subspace-measure}
设 $X$ 为 LCH 空间，$Y$ 为 $X$ 的闭子集（赋予相对拓扑后也是 LCH 空间），$\mu$ 为 $Y$ 上的 Radon 测度。则 $I(f)=\int(f|_Y)\,\mathrm d\mu$ 是 $C_c(X)$ 上的正线性泛函，并且它在 $X$ 上诱导的 Radon 测度为 $\nu(E)=\mu(E\cap Y)$。
\end{exercise}
% SUPPLEMENT: id=v14-ch07-ex01-solution; parent=ch07-ex01; kind=solution+motivation+hypotheses; origin=independent-derivation
\begin{solution}
\textbf{闭子空间与延到全空间的测度（重构新增）。}Hausdorff 性遗传到 $Y$。若 $y\in Y$，取 $X$ 中 $y$ 的紧邻域 $K$；由于 $Y$ 闭，$K\cap Y$ 紧且包含 $y$ 的相对开邻域，故 $Y$ 局部紧。对 $f\in C_c(X)$，$f|_Y$ 连续，其在 $Y$ 中的支集包含于 $\operatorname{supp}(f)\cap Y$，后者紧，所以 $f|_Y\in C_c(Y)$。由紧集有限性它关于 $\mu$ 可积，积分线性并保持非负性，得所述 $I$。

对 Borel 集 $E\subseteq X$，交集 $E\cap Y$ 是 $Y$ 中的 Borel 集。因此 $\nu(E)=\mu(E\cap Y)$ 定义非负可数可加测度；对紧 $K\subseteq X$，$\nu(K)<\infty$。为证外正则，给定 Borel 集 $E$，若 $\nu(E)=\infty$ 则自动成立；若有限，取 $Y$ 中的相对开集 $V\supseteq E\cap Y$ 使 $\mu(V)<\nu(E)+\varepsilon$。写 $V=O\cap Y$，其中 $O$ 在 $X$ 中开。因为 $Y$ 闭，
\begin{align*}
W=O\cup(X\setminus Y)
\end{align*}
在 $X$ 中开，包含 $E$，而 $W\cap Y=V$，所以 $\nu(W)<\nu(E)+\varepsilon$。这证明外正则性。

若 $U\subseteq X$ 开，则 $U\cap Y$ 相对开。由 $\mu$ 的内正则性，$\mu(U\cap Y)$ 为其紧子集 $K\subseteq U\cap Y$ 的测度上确界。每个这样的 $K$ 经包含映射仍是 $X$ 中的紧集，且 $\nu(K)=\mu(K)$，故 $\nu$ 在 $U$ 上内正则。因此 $\nu$ 是 Radon 测度。对示性函数，积分恒等式 $\int_X\chi_E\,\mathrm d\nu=\int_Y\chi_E|_Y\,\mathrm d\mu$ 就是定义；由简单函数、非负单调逼近及实虚部分分解，得到所有可积函数的同一恒等式。特别地 $I(f)=\int_Xf\,\mathrm d\nu$，由\cref{theorem:ch07-riesz-positive}的唯一性，$\nu$ 正是诱导测度。
\end{solution}
\begin{remark}
\textbf{动机与逐条件说明（重构新增）。}这个构造把 $Y$ 上的质量放入 $X$，在 $X\setminus Y$ 上不给质量。$Y$ 闭既保证紧支函数限制后仍紧支，也使外逼近时可以加入开集 $X\setminus Y$；若随意删除闭性，这两个步骤都不再自动成立。Radon 三条件分别在紧集、开集和 Borel 集上验证，不能只由可数可加性就宣称诱导测度为 Radon。
\end{remark}
% END SUPPLEMENT: id=v14-ch07-ex01-solution

% SOURCE p=230; printed=215; object=ch07-ex02; number=2; region=lower
\begin{exercise}\label{exercise:ch07-measure-support}
设 $\mu$ 为 $X$ 上的 Radon 测度。
\begin{enumerate}[label=\alph*.,ref=(\alph*)]
\item\label{item:ch07-support-null-open-union} 令 $N$ 为所有满足 $\mu(U)=0$ 的开集 $U\subseteq X$ 的并。则 $N$ 开且 $\mu(N)=0$。
\begin{definition*}[测度的支集]\phantomsection\label{definition:ch07-measure-support}
$N$ 的补集称为 $\mu$ 的\textbf{支集}，记为 $\operatorname{supp}(\mu)$。
\end{definition*}
\item\label{item:ch07-support-test-functions} $x\in\operatorname{supp}(\mu)$ 当且仅当，对每个满足 $f(x)>0$ 的 $f\in C_c(X,[0,1])$，都有 $\int f\,\mathrm d\mu>0$。
\end{enumerate}
\end{exercise}
% SUPPLEMENT: id=v14-ch07-ex02-solution; parent=ch07-ex02; kind=solution+motivation+hypotheses; origin=independent-derivation
\begin{solution}
\textbf{任意零开集并与连续测试（重构新增）。}(a) 任意开集的并仍开，故 $N$ 开。若 $K\subseteq N$ 紧，则定义 $N$ 的零测度开集族覆盖 $K$，从中取有限子覆盖 $U_1,\ldots,U_m$，得 $\mu(K)\leqslant\sum_{j=1}^m\mu(U_j)=0$。由 $N$ 上的内正则性，
\begin{align*}
\mu(N)=\sup\{\mu(K):K\subseteq N,\ K\text{ 紧}\}=0.
\end{align*}
因此支集为闭集，且其补集是零集。

(b) 若 $x\in\operatorname{supp}(\mu)$，则任何含 $x$ 的开集都不可能测度为零。给定 $f$ 如题，取 $0<c<f(x)$。连续性使 $U=\{y:f(y)>c\}$ 为含 $x$ 的开集，故 $\mu(U)>0$，且 $\int f\,\mathrm d\mu\geqslant c\mu(U)>0$。若 $x\notin\operatorname{supp}(\mu)$，则 $x\in N$；由局部 Urysohn 引理，存在 $f\in C_c(X,[0,1])$，使 $f(x)=1$ 且 $\operatorname{supp}(f)\subseteq N$。于是 $\int f\,\mathrm d\mu=0$，与题设的全称测试条件矛盾。这证明反向。
\end{solution}
\begin{remark}
\textbf{动机与逐条件说明（重构新增）。}支集保留每个邻域都有质量的点；连续非负峰函数便能在这样的点附近检测正质量。任意多个零测度集合的并一般并非零测度，这里借助紧集的有限子覆盖及\emph{开集内正则性}才得出 $\mu(N)=0$。LCH 条件用于在 $N$ 内构造峰函数，紧集有限性使这些函数的积分有限。测试条件必须要求 $f(x)>0$，否则零函数会立即破坏结论；$f\geqslant0$ 排除积分中的正负抵消，而上界一是方便的归一化。支集中的点未必有正的单点质量，例如 Lebesgue 测度的每个点都在支集内，却没有单点质量。
\end{remark}
% END SUPPLEMENT: id=v14-ch07-ex02-solution

% SOURCE p=230; printed=215; object=ch07-ex03; number=3; region=lower
\begin{exercise}\label{exercise:ch07-discrete-compactification-support}
设 $X$ 为一个离散拓扑集合的单点紧化。若 $\mu$ 是 $X$ 上的 Radon 测度，则 $\operatorname{supp}(\mu)$（见\cref{exercise:ch07-measure-support}）可数。
\end{exercise}
% SUPPLEMENT: id=v14-ch07-ex03-solution; parent=ch07-ex03; kind=solution+motivation+hypotheses; origin=independent-derivation
\begin{solution}
\textbf{有限总质量控制孤立点（重构新增）。}写 $X=D\cup\{\infty\}$，其中 $D$ 离散。$X$ 紧，故 $M=\mu(X)<\infty$。对 $x\in D$，单点 $\{x\}$ 是开集；因此 $x$ 在支集中当且仅当 $\mu(\{x\})>0$：零质量单点本身是零开邻域，而正质量单点使任何邻域都有正质量。

令 $D_n=\{x\in D:\mu(\{x\})\geqslant1/n\}$。若其中有 $m$ 个不同点，有限可加性给出 $m/n\leqslant M$，故 $D_n$ 必为有限集。所有正质量孤立点组成 $\bigcup_{n=1}^\infty D_n$，是可数集。再至多添上一个点 $\infty$，支集仍可数。
\end{solution}
\begin{remark}
\textbf{动机与逐条件说明（重构新增）。}离散性使原集合中的点可由单点开邻域测试，紧化加上 Radon 的紧集有限性则给出有限总质量，禁止不可数个正质量原子。新增点 $\infty$ 即使没有单点质量，也可能属于支集，因此证明只说“至多添上一个点”，未把全部支集都误认为原子。去掉紧化后的有限总质量约束，任意不可数离散集上的计数测度就是 Radon 测度，支集却为整个不可数空间：紧集只有有限集，所有集合均开，三项 Radon 条件直接成立。
\end{remark}
% END SUPPLEMENT: id=v14-ch07-ex03-solution

% SOURCE p=230,231; printed=215-216; object=ch07-ex04; number=4; region=bottom-top
\begin{exercise}\label{exercise:ch07-compact-gdelta}
设 $X$ 为 LCH 空间。
\begin{enumerate}[label=\alph*.,ref=(\alph*)]
\item\label{item:ch07-level-compact-gdelta} 若 $f\in C_c(X,[0,\infty))$，则对每个 $a>0$，$f^{-1}([a,\infty))$ 是紧 $G_\delta$ 集。
\item\label{item:ch07-compact-gdelta-level-one} 若 $K\subseteq X$ 为紧 $G_\delta$ 集，则存在 $f\in C_c(X,[0,1])$ 使 $K=f^{-1}(\{1\})$。
\item\label{item:ch07-baire-generated-compact-gdelta} Baire 集的 $\sigma$ 代数 $\mathscr B_X^0$ 正是由紧 $G_\delta$ 集生成的 $\sigma$ 代数。
\end{enumerate}
\end{exercise}
% SUPPLEMENT: id=v14-ch07-ex04-solution; parent=ch07-ex04; kind=solution+motivation+hypotheses; origin=independent-derivation
\begin{solution}
\textbf{水平集刻画与生成元（重构新增）。}(a) 集合 $K=\{f\geqslant a\}$ 闭，且因 $a>0$ 包含于紧集 $\operatorname{supp}(f)$，故紧。同时
\begin{align*}
K=\bigcap_{n=1}^\infty\{x:f(x)>a-1/n\},
\end{align*}
右端为可数个开集之交，故为 $G_\delta$ 集。

(b) 若 $K=\varnothing$，取 $f=0$。否则写 $K=\bigcap_{n=1}^\infty U_n$，其中每个 $U_n$ 开。由局部 Urysohn 引理，取 $\phi_0\in C_c(X,[0,1])$ 使其在 $K$ 上为一，并对每个 $n\geqslant1$ 取 $\phi_n\in C_c(X,[0,1])$，使 $\phi_n=1$ 于 $K$ 且 $\operatorname{supp}(\phi_n)\subseteq U_n$。令
\begin{align*}
f=\phi_0\sum_{n=1}^\infty2^{-n}\phi_n.
\end{align*}
由于每项模不超过 $2^{-n}$，级数一致收敛，故其和连续；乘以 $\phi_0$ 后支集包含于紧集 $\operatorname{supp}(\phi_0)$。因此 $f\in C_c(X,[0,1])$，且在 $K$ 上等于一。若 $x\notin K$，存在 $n$ 使 $x\notin U_n$，从而 $\phi_n(x)=0$，其余项至多为其权重；于是 $\sum_j2^{-j}\phi_j(x)\leqslant1-2^{-n}<1$，故 $f(x)<1$。这证明 $f^{-1}(\{1\})=K$。

(c) 记紧 $G_\delta$ 集生成的 $\sigma$ 代数为 $\mathscr G$。由 (b)，每个生成元为某个 $C_c$ 函数的闭集原像，故为 Baire 集，所以 $\mathscr G\subseteq\mathscr B_X^0$。反过来，若 $h\in C_c(X,[0,\infty))$，则由 (a)，$\{h\geqslant a\}\in\mathscr G$ 对所有 $a>0$ 成立。对 $t\geqslant0$，
\begin{align*}
\{h>t\}=\bigcup_{\substack{r\in\mathbb Q\\r>t}}\{h\geqslant r\};
\end{align*}
而 $t<0$ 时此集合为 $X$，故 $h$ 关于 $\mathscr G$ 可测。任意复值 $f\in C_c(X)$ 的实部、虚部各自的正负部分都属于上述非负函数类，所以 $f$ 也关于 $\mathscr G$ 可测。由 $\mathscr B_X^0$ 的最小性，$\mathscr B_X^0\subseteq\mathscr G$，两者相等。
\end{solution}
\begin{remark}
\textbf{动机与逐条件说明（重构新增）。}正高度的闭水平集是紧支连续函数最直接的紧集合数据。(a) 中 $a>0$ 必须保留：若 $a=0$ 且 $f\geqslant0$，原像为整个 $X$，在非紧空间中不紧。(b) 的紧性提供每个截断函数，而可数交条件允许用全部正权重的级数编码每一个开集限制；另乘 $\phi_0$ 是为了把级数的支集统一限制在一个紧集内，不能仅凭各项分别紧支就断言其和紧支。(c) 的生成结论只涉及 $\sigma$ 代数，而不要求每个 Baire 集本身都是紧 $G_\delta$ 集。
\end{remark}
% END SUPPLEMENT: id=v14-ch07-ex04-solution

% SOURCE p=231; printed=216; object=ch07-ex05; number=5; region=upper
\begin{exercise}\label{exercise:ch07-second-countable-baire}
设 $X$ 为第二可数的 LCH 空间。
\begin{enumerate}[label=\alph*.,ref=(\alph*)]
\item\label{item:ch07-second-countable-compact-gdelta} $X$ 的每个紧子集都是 $G_\delta$ 集。
\item\label{item:ch07-second-countable-baire-borel} $\mathscr B_X^0=\mathscr B_X$。
\end{enumerate}
\end{exercise}
% SUPPLEMENT: id=v14-ch07-ex05-solution; parent=ch07-ex05; kind=solution+motivation+hypotheses; origin=independent-derivation
\begin{solution}
\textbf{把局部选择化成可数操作（重构新增）。}(a) 取可数开基 $(B_n)$。对紧集 $K$ 之外的每个点 $x$，局部紧 Hausdorff 空间的缩小邻域性质给出开邻域 $V_x$，使 $\overline{V_x}\cap K=\varnothing$；具体可在开集 $X\setminus K$ 中取紧邻域，再取其内部。选一个基元 $B_n$ 使 $x\in B_n\subseteq V_x$，于是 $\overline{B_n}\cap K=\varnothing$。令 $J=\{n:\overline{B_n}\cap K=\varnothing\}$，则
\begin{align*}
X\setminus K=\bigcup_{n\in J}\overline{B_n},
\qquad K=\bigcap_{n\in J}(X\setminus\overline{B_n}).
\end{align*}
右端为可数个开集之交。若 $J$ 为空，交集按约定为 $X$，也可将它写为 $\bigcap_{n=1}^\infty X$，故仍为 $G_\delta$ 集。

(b) 给定开集 $U$，对每个 $x\in U$，选开邻域 $V_x$ 使 $\overline{V_x}$ 紧且包含于 $U$。第二可数空间的每个子空间都第二可数，从而 $U$ 是 Lindelöf 空间；因此 $(V_x)_{x\in U}$ 有可数子覆盖 $(V_{x_n})$。于是
\begin{align*}
U=\bigcup_{n=1}^\infty\overline{V_{x_n}}.
\end{align*}
由 (a)，每个闭包都是紧 $G_\delta$ 集；由\subcref{exercise:ch07-compact-gdelta}{item:ch07-baire-generated-compact-gdelta}，它们均为 Baire 集。所以所有开集都是 Baire 集，得到 $\mathscr B_X\subseteq\mathscr B_X^0$。反向包含已由连续函数的 Borel 可测性给出。
\end{solution}
\begin{remark}
\textbf{动机与逐条件说明（重构新增）。}第二可数性把逐点的邻域选择压缩到可数个，从而让补集、可数并交这些 $\sigma$ 代数操作足以恢复全部开集。局部紧性保证用于覆盖的闭包紧，Hausdorff 性保证紧集闭，并支持在补集内缩小邻域。(a) 实际使用的是可数基加这种分离性质；(b) 还使用开集内的预紧邻域。若只假设全空间 Lindelöf，不能直接断言每个开子空间也 Lindelöf，因此这里明确利用遗传的第二可数性。
\end{remark}
% END SUPPLEMENT: id=v14-ch07-ex05-solution

% SOURCE p=231; printed=216; object=ch07-ex06; number=6; region=upper
\begin{exercise}\label{exercise:ch07-baire-borel-different}
设 $X$ 是赋予离散拓扑的不可数集，或这种集合的单点紧化。则 $\mathscr B_X\ne\mathscr B_X^0$。
\end{exercise}
% SUPPLEMENT: id=v14-ch07-ex06-solution; parent=ch07-ex06; kind=solution+motivation+hypotheses; origin=independent-derivation
\begin{solution}
\textbf{两种空间分别计算（重构新增）。}先设 $X=D$ 为不可数离散集。其紧集恰为有限集：若某集合无限，用单点开集覆盖便无有限子覆盖；反之有限集总紧。所有有限集都开，因此也是 $G_\delta$ 集。由\subcref{exercise:ch07-compact-gdelta}{item:ch07-baire-generated-compact-gdelta}，$\mathscr B_D^0$ 是有限集生成的 $\sigma$ 代数，即可数集与余可数集所成的 $\sigma$ 代数。确实，这一族是包含单点集的 $\sigma$ 代数，而每个可数集是单点的可数并，其补集也在生成的 $\sigma$ 代数中。离散拓扑的 Borel $\sigma$ 代数则为 $\mathscr P(D)$。将不可数集 $D$ 分为两个不可数子集，任取其中一个，便是 Borel 集而非 Baire 集（这里沿用本书的选择公理背景）。

再设 $X=D\cup\{\infty\}$ 为单点紧化。含 $\infty$ 的开邻域的补集是 $D$ 的紧集，因而有限。不含 $\infty$ 的紧子集是离散空间 $D$ 的紧集，故有限。若紧 $G_\delta$ 集 $K$ 含 $\infty$，写 $K=\bigcap_n U_n$，则每个 $U_n$ 都含 $\infty$，从而
\begin{align*}
X\setminus K=\bigcup_n(X\setminus U_n)
\end{align*}
是 $D$ 中的可数集。因此所有紧 $G_\delta$ 集都属于
\begin{align*}
\mathscr A=\{A\subseteq D:A\text{ 可数}\}
\ \cup\ \{X\setminus A:A\subseteq D,\ A\text{ 可数}\}.
\end{align*}
这个族是 $\sigma$ 代数：取补集互换两类；可数并若全属第一类仍可数，若至少有一个第二类，其补集包含于某个可数的 $D$ 子集，也属于第一类。每个 $x\in D$ 的单点集又是紧开集，故为紧 $G_\delta$ 集；这些单点生成 $\mathscr A$。由上一题的生成结论，$\mathscr B_X^0=\mathscr A$。但 $D$ 在 $X$ 中开，因此为 Borel 集；它不可数且不含 $\infty$，不属于 $\mathscr A$。故此时两者也不相等。
\end{solution}
\begin{remark}
\textbf{动机与逐条件说明（重构新增）。}连续紧支函数通过可数次测度操作能区分的集合，少于任意开集族所包含的信息。不可数性使这种差距出现；在可数离散集上，单点的可数并已经给出全部子集。第二种情形虽已把空间变紧，仍不能补回缺少可数基带来的问题。其 Baire $\sigma$ 代数必须区分新增点 $\infty$：所有不含它的 Baire 集都可数，不能直接套用 $X$ 上未经区分的“可数或余可数”说法。
\end{remark}
% END SUPPLEMENT: id=v14-ch07-ex06-solution
\end{document}
